% Transcription — fonds Alexandre Grothendieck, shelfmark 39, batch 1 (pages 1-19).
% Established from the facsimile archives/batches/39/batch-01.pdf (a mirror of
% https://grothendieck.umontpellier.fr/39.pdf), per the skill
% transcribe-grothendieck. The folder is nineteen pages and this is its only
% batch.
%
% Sixteen of the nineteen leaves are transcribed. Pages 1, 8 and 17 are blank
% cream separator sheets carrying nothing in his hand; they take no \page
% marker. Nothing in the folder carries a title of his except page 18's
% « Th. de préparation "formel" »; the other two section headings are the
% edition's and say so. He does not paginate.
%
% Three runs, in blue ink throughout. Pages 2-7: k a field, S_n(k) =
% k[[t_1..t_n]] as a cosimplicial ring, and a cosimplicial subring A_* subject
% to a list of conditions — a) stability under the symmetric group, adding a
% variable, repeating a variable; b) unital subrings; c) local with local
% inclusion; d) t_1 in A_1; e) closure under composition of « A-morphisms »
% X_m -> X_n; f) inverse of an A-automorphism is an A-morphism, implying f')
% the A_n are henselian — with the smallest A_* for each set (Im Z, prime
% field, P[t], its localisation, its henselisation). Then, in characteristic
% 0, g), g') exponentials of formal vector fields and solutions of formal ODEs
% are A-morphisms; the Hausdorff (Campbell-Hausdorff) group law and Ado reduce
% « formal Lie groups over k_0 = A_0 are A-groups » to linear groups; A_1
% contains log(1+t), exp, cos, sin, tg; A-foliations and h): a completely
% integrable A-Pfaff system gives an A-foliation. Page 7 is a cancelled
% restart of the end of page 6 and is transcribed because it reads.
% Pages 9-16: a second redaction over an arbitrary commutative ring k, with
% X_* = (Spf S_n) an augmented simplicial formal scheme: a) identities, b)
% composition, c) « difference » and « product » are A_*-morphisms (the
% elementary conditions); the smallest such A_* is Lambda[t_1..t_n]; the
% category of A_*-formal schemes, the functor from free k_0-modules; d) the
% henselian condition (a formal isomorphism that is an A_*-morphism has
% A_*-inverse), with Cor. 1 (A_n cap m_n in rad A_n — and his own objection
% in the margin that the substitution t_{n+1} = 1 is not licensed), Cor. 2
% and Cor. 3 (implicit function / Hensel); sub-A_*-formal schemes and
% uniqueness of induced structure; a differential condition (coefficients of
% the expansion in y lie in A_n) giving A_*-structures on P^n_{X/k}, normal
% invariants, Diff^n; quotient A_*-structures along a submersion and f) the
% « condition intégrale »: a formal equivalence relation that is an
% A_*-substructure has an A_*-quotient, heuristically the integrability of a
% completely integrable Pfaff system by an A_*-morphism. Pages 18-19: the
% formal Weierstrass preparation theorem over a linearly topologised
% separated complete ring, proved by reduction to A discrete and a graded
% bijectivity argument for (b, Q) |-> sum b_i Z^i + QF; a first attempt via
% (i) Nakayama, (ii) regularity, (iii) the monic polynomial case is boxed and
% cancelled.
%
% Reading hazards fixed once here. His A_* is written with a small star-like
% subscript that reads as « A_x » or « A_n »; it is set A_* wherever the
% cosimplicial family is meant and A_n where an index is. The script O of
% O_{A_*}(X) and the script P of P^n_{X/k} are set \mathcal. « hom. »,
% « isom. », « morph. », « syst. », « compl. », « op. diff. », « rel.
% d'équiv. », « corr. biunivoque », « gén. », « cqfd » are his abbreviations
% and kept. His underlining is set in bold. The letters of the conditions do
% not form one sequence across the second redaction: page 11 calls the
% henselian condition d), page 14 opens a differential condition also
% lettered d), pages 13 and 15 invoke « e) », and page 15 continues with f);
% the letters are transcribed as written and the collision noted on page 14.
% On pages 18-19 the dot marks reduction mod m (a dot over a_i, F, u).
%
% Readings flagged and not settled: the marginal notes of pages 3 (bottom),
% 4, 10 (the cancelled boxed note), 12 and 14, which are small and fast;
% « unifère » on page 2; « le déterminant » on page 3; « intervention » on
% page 4; « (d) » on page 13; the reduction step opening the proof on page
% 18; the value of alpha on page 19. Page 2 writes 1/f in S_n where the
% argument wants A_n; it is flagged, not repaired.
%
% Pass: Opus 5 (claude-opus-5), 2026-09-13 — first pass, unchecked against the pages by a human.
\documentclass[11pt,a4paper]{article}
% Le préambule commun à toutes les transcriptions du fonds Grothendieck.
%
% Il tient en une idée : **l'incertitude se compose, elle ne se résout pas.**
% Une transcription qui lisse un mot douteux a détruit la seule chose pour
% laquelle on la faisait — un lecteur doit pouvoir savoir, à chaque endroit,
% ce qui était sur le feuillet et ce qui a été ajouté. Les six macros qui
% suivent sont donc l'appareil critique, et non de la décoration.
%
% Les mêmes commandes sont reconnues par `scripts/render.mjs`, qui produit la
% vue de lecture du site. Une macro ajoutée ici doit l'être aussi là-bas,
% faute de quoi le rendu échouera bruyamment — ce qui est voulu : un rendu
% silencieusement faux est pire qu'un rendu absent.

\ProvidesPackage{grothendieck}[2026/08/08 Transcriptions du fonds Grothendieck]

\usepackage{amsmath,amssymb,amsthm}
\usepackage{mathrsfs}
\usepackage{stmaryrd}
% Les diagrammes commutatifs sont partout dans ce fonds : c'est la langue
% même de la géométrie algébrique de Grothendieck, et une transcription qui
% les rendrait en prose ne transcrirait rien.
\usepackage{tikz-cd}
% Français par défaut — c'est la langue des feuillets — mais l'anglais est
% chargé aussi : la traduction et le résumé sont des documents anglais, et la
% typographie française (espaces avant les deux-points) y serait une faute.
% Ces documents basculent par \selectlanguage{english} en tête de corps.
\usepackage[english,french]{babel}
\usepackage{fontspec}
\usepackage{geometry}
\geometry{a4paper,margin=25mm,marginparwidth=28mm}
\usepackage{marginnote}
\usepackage{xcolor}
% ulem, not soul. `soul`'s \st and \ul are built on a character-by-character
% scan that XeLaTeX's font machinery breaks: the document fails with an
% undefined control sequence rather than degrading. ulem does the same two jobs
% and is Unicode-engine safe. `normalem` stops it hijacking \emph, which the
% transcriptions use for Grothendieck's own emphasis.
\usepackage[normalem]{ulem}
\usepackage{hyperref}
\hypersetup{colorlinks=true,linkcolor=black,urlcolor=black,pdfborder={0 0 0}}

% --- Métadonnées du lot -----------------------------------------------------
% Reprises telles quelles par le script de rendu, qui les affiche en tête de
% la vue de lecture. Elles fixent ce que le fichier prétend couvrir : sans
% elles, une transcription flotte sans point d'attache dans un fonds de seize
% mille feuillets.

\newcommand{\folder}[1]{\gdef\grothFolder{#1}}
\newcommand{\batch}[1]{\gdef\grothBatch{#1}}
\newcommand{\leaves}[2]{\gdef\grothFirst{#1}\gdef\grothLast{#2}}
\newcommand{\foldertitle}[1]{\gdef\grothFoldertitle{#1}}
\gdef\grothFolder{?}\gdef\grothBatch{?}\gdef\grothFirst{?}\gdef\grothLast{?}\gdef\grothFoldertitle{}

% --- Le résumé --------------------------------------------------------------
%
% Il ouvre la lecture modernisée, sous le titre « Résumé », et il est composé
% en petit. Ce n'est pas un ornement : c'est le seul passage du document qui
% ne soit pas des mathématiques, et un lecteur doit pouvoir le reconnaître
% avant de l'avoir lu — puis le sauter s'il connaît déjà le sujet, ou s'y
% arrêter s'il ne le connaît pas. Le filet qui le suit marque la reprise du
% corps.

\newenvironment{resume}
  {\small\setlength{\parskip}{0.45em}}
  {\par\medskip\hrule\medskip}

% --- Mots-clés ---------------------------------------------------------------
%
% En anglais, en clôture du résumé de la lecture modernisée : le vocabulaire
% moderne sous lequel chercher ce que les feuillets construisent. Le manifeste
% du site (`npm run manifest`) les extrait pour étiqueter la cote entière —
% cette ligne est la source unique des tags, il n'y a pas de fichier à côté.
\newcommand{\keywords}[1]{\par\smallskip\noindent
  {\footnotesize\textsc{Keywords} --- \textit{#1}}\par}

% --- L'appareil critique ----------------------------------------------------

% Le numéro de feuillet, en marge. C'est l'ancre qui permet de régler un
% désaccord sur une formule en regardant *un* feuillet plutôt qu'une chemise
% de six cents. Le site s'en sert aussi pour tourner les pages du fac-similé
% quand on fait défiler la transcription.
\newcommand{\leaf}[1]{%
  \marginnote{\footnotesize\color{gray!70}\textsf{#1}}%
  \label{leaf:#1}\ignorespaces}

% Illisible : on ne devine pas. Un mot inventé qui se lit comme les autres est
% le pire résultat possible.
\newcommand{\ill}{\textcolor{red!60!black}{[\,\ldots\,]}}

% Lecture douteuse : le mot est donné, et signalé comme incertain.
\newcommand{\uncertain}[1]{\uline{#1}}

% Ajout de l'éditeur — une lettre manquante, un mot sous-entendu.
\newcommand{\add}[1]{\textcolor{blue!50!black}{[#1]}}

% Biffé par Grothendieck lui-même. Ce qu'il a rayé fait partie du document :
% une rature dit souvent où la pensée a changé de direction.
\newcommand{\struck}[1]{\sout{#1}}

% Note du transcripteur, sur l'état matériel du feuillet.
\newcommand{\note}[1]{\par\smallskip{\small\color{gray!60!black}\textit{[#1]}}\par\smallskip}

% Note marginale de Grothendieck, distincte du corps du texte.
\newcommand{\marginal}[1]{\par\smallskip\noindent\hspace{1em}%
  {\small\color{gray!60!black}#1}\par\smallskip}

% --- Titre ------------------------------------------------------------------

\AtBeginDocument{%
  \begin{center}
    {\large\bfseries Fonds Alexandre Grothendieck --- cote n\textsuperscript{o}~\grothFolder}\\[2pt]
    {\itshape\grothFoldertitle}\\[4pt]
    {\small feuillets \grothFirst--\grothLast\ (lot \grothBatch)}\\[2pt]
    {\footnotesize\color{gray!60!black}Transcription établie d'après le fac-similé numérisé
     par l'Université de Montpellier.\\
     Les lectures incertaines sont soulignées, les illisibles notées [\,\ldots\,],
     les ajouts entre crochets.}
  \end{center}
  \vspace{1em}}

\endinput


\folder{39}
\batch{1}
\pages{1}{19}
\dating{[vers 1960-1970]}
\watermark{Édition de démonstration}
\foldertitle{Fonctions et structures « élémentaires » [Schémas formels, théorème de préparation] : notes manuscrites (s.d.)}

\begin{document}

\section{Sous-anneaux cosimpliciaux de $k[[t_1, \ldots, t_n]]$, $k$ corps}

\note{ce titre est de l'édition : aucune page du dossier ne porte de titre
avant la page 18. Les pages 1, 8 et 17 sont des feuillets intercalaires
vierges}

\page{2}

$k$ un \struck{corps \ill{} \ill{}} corps

$S_n(k) = k[[t_1, \ldots, t_n]]$ \quad ($n \geq 0$ ; $S_0(k) = k$) \quad
\uncertain{forment}, \uncertain{pour $n$ variable}, un anneau cosimplicial

$A_n \subset S_n(k)$ \quad sous-anneau \uncertain{cosimplicial} \uncertain{t.q.}

\begin{itemize}
  \item[(b)] $\forall n \geq 0$, $A_n$ est un sous-anneau \uncertain{unifère} de $S_n$
  \item[(a)] on a
  \item[a$_1$)] $A_n$ est stable par $\mathfrak{S}_n$ opérant sur $S_n$
  \item[a$_2$)] si $f(t_1, \ldots, t_n) \in A_n$, alors
    $\underbrace{f(t_1, \ldots, t_n)}$ $\in A_{n+1}$
    \marginal{$t_{n+1}$ variable qui ne figure pas}
  \item[a$_3$)] si $f(t_1, \ldots, t_n) \in A_n$, alors
    $f(t_1, t_2, \ldots, \underbrace{t_{n-1}, t_{n-1}}_{\text{répété}}) \in A_{n-1}$
\end{itemize}

\note{une flèche ramène (b) sous (a) : l'ordre voulu est a), b)}

Le plus petit $A_* = (A_n)_{n \in \mathbf{N}}$ \uncertain{vérifiant a) b)}
\uncertain{famille} est donné par \struck{$A_* = \ill{}$}
\[
  A_n = \mathrm{Im}(\mathbf{Z} \to S_n)
\]
dans l'\struck{objet} \add{anneau} cosimplicial constant de valeur
$\mathrm{Im}(\mathbf{Z} \to k)$ (qui est $\mathbf{Z}$ ou un $\mathbf{Z}/p\mathbf{Z}$,
\struck{\ill{}} \uncertain{$p$ nombre premier} $= \mathrm{car}.\, k$)\note{« anneau »
est écrit au-dessus de « objet » biffé}

On suppose de plus que

\begin{itemize}
  \item[c)] si $f \in A_n$ et $f \notin \mathrm{rad}\, S_n$, i.e. $f$ inv. dans
    $S_n$, alors $\frac{1}{f} \in S_n$\note{lire sans doute $\frac{1}{f} \in A_n$}
  \item[] $\Updownarrow$
  \item[c$'$)] $A_n$ est un anneau local, et l'inclusion $A_n \to S_n$ est un
    hom. local
\end{itemize}

Le plus petit $A_*$ satisfaisant a) b) c) est formé des
\[
  A_n = \mathrm{Im}(P \to S_n)
\]
où $P$ est le sous-corps premier de $k$ ($\mathbf{Q}$ en car. $0$,
$\mathbf{F}_p$ en car. $p > 0$).

\begin{itemize}
  \item[d)] $t_1 \in A_1 \subset S[[t_1]]$
\end{itemize}

Le plus petit $A_*$ satisfaisant a) b) d) \struck{c)} est formé des
\[
  A_n = P[t_1, \ldots, t_n]
\]
\struck{\ill{}}
le plus petit $A$ satisfaisant a) b) c) d) est formé des
\[
  A_n = P[t_1, \ldots, t_n]_{\mathfrak{m}_n}, \qquad
  \mathfrak{m}_n = \sum t_i\, P[t_1, \ldots, t_n].
\]

\struck{e) \ill{}}

Un hom. $S_n \to S_m$ (ou $\mathrm{Spf}(S_m) \to \mathrm{Spf}(S_n)$)\note{sous
les deux $\mathrm{Spf}$ il écrit $X_m$ et $X_n$ : $X_n = \mathrm{Spf}(S_n)$ ; son
$X$ bouclé se lit d'abord comme un $\mathcal{D}$} donné par
$f_1, \ldots, f_n \in S_m$ est dit un $A$-hom. si $f_1, \ldots, f_n \in A_m$. Il
résulte de b) et d) que l'identité $X_n \to X_n$, plus généralement les
applications \struck{et \ill{}} $X_n \to X_m$ \ill{}
\uncertain{sont des} $A$-hom. ($t_i \in A_n$ pour $1 \leq i \leq n$). L'analogue
de b)

\page{3}

est la condition

\begin{itemize}
  \item[e)] Le composé de deux ($A_*$-)morphismes $X_m \to X_n \to X_p$ est un
    $A_*$-morphisme.
\end{itemize}

\marginal{Moyennant a) à e), \uncertain{i.e.}
$\mathrm{Hom}_A(X_m, X_n) \simeq \mathrm{Hom}_{A_0\text{-alg}}(A_n, A_m)$}

\textbf{NB} Moyennant b) et e), la condition a) équivaut au \struck{fait} que
\[
  \begin{cases}
    \text{a}_1)\ 1 \in A_0 \\
    \text{a}_2)\ \text{le morphisme « différence » } X_1 \times X_1 \to X_1
      \text{ est } A_*\text{-morph.} \\
    \text{a}_3)\ \text{le morphisme « produit » } X_1 \times X_1 \to X_1
      \text{ est } A_*\text{-morph.}
  \end{cases}
\]

\textbf{NB} Les $A_n = k[t_1, \ldots, t_n]$ ou $A_n = k[t_1, \ldots, t_n]_{\mathfrak{m}_n}$
satisfont e)

L'analogue de c) est la condition

\begin{itemize}
  \item[f)] Si $u : X_n \to X_n$ est un $A$-morphisme et un isomorphisme,
    \struck{alors} i.e. le \uncertain{déterminant} \struck{\ill{}} de la matrice
    jacobienne de $u$ est inversible, alors $u^{-1}$ est un $A$-morphisme.
\end{itemize}

Cette condition implique \struck{(et \ill{} équivalente} \struck{à ??}

\begin{itemize}
  \item[f$'$)] Les $A_n$ sont des anneaux henséliens
\end{itemize}

\marginal{f$'$) implique-t-il f) (moyennant a) à e)) ? il \uncertain{suffit} par
\uncertain{récurrence}}

[En fait, si f) est satisfait pour un $n$, (f$'$) est satisfait pour $n - 1$, i.e.
$A_{n-1}$ est hensélien] \struck{D'ailleurs, si \struck{$X_n \simeq X_n$} $u$ est
donné par $f_1, \ldots, f_n \in B_n$, alors $u^{-1}$ est donné par $g_1, \ldots,
g_n$ appartenant à la clôture hensélienne de $(k[f_1, \ldots, f_n]_{\mathfrak{n}}$,
où $\mathfrak{n} = (k[f_1, \ldots, f_n]$}\note{le passage « D'ailleurs … » est barré
de longs traits obliques}

\textbf{NB} Le plus petit $A_*$ satisfaisant a) à d) et f$'$) est formé des
\[
  A_n = P\{t_1, \ldots, t_n\} \qquad
  \bigl(\text{clôture hensélienne de } P[t_1, \ldots, t_n]_{\mathfrak{m}_n}\bigr)
\]
et il satisfait e) et f).

\struck{\ill{} changement de variables for\ill{}}

Moyennant a) b) c) d) e) \struck{on peut} \uncertain{définir} la catégorie des
$A$-schémas formels, et des Modules dessus (\uncertain{faisceaux} \ill{}), et le
calcul différentiel \uncertain{fini} dessus (opérateurs différentiels, parties
principales, etc.\ \ldots)

\marginal{\ill{} (\uncertain{6 des propriétés}) \ill{} \uncertain{produits fibrés}
\uncertain{transversaux} \ill{}}

\page{4}

Nous supposons maintenant $k$ de car.~$0$, de sorte qu'on a la notion
d'exponentielle d'un champ de vecteurs formel $\xi$ sur \struck{$X$} un schéma
formel lisse $X$,
\[
  (t, x) \mapsto (\exp t\xi) . x \;:\; \underline{X_1} \times X \to X
\]
\note{sous $X_1$, en petit : « \uncertain{le temps} »}

\marginal{Indiquer \ill{} d'équations différentielles \ill{}}

\begin{itemize}
  \item[g)] Si $X$ est un $A$-schéma formel (lisse) et $\xi$ un $A$-champ de
    vecteurs, alors le morphisme $X_1 \times X \to X$ est un $A$-morphisme.
\end{itemize}

(Bien entendu, il suffit de le faire \add{dans le cas} des $X$ de la forme
$X_n$\ldots).

Ceci indique (vu l'\uncertain{intervention} de la variable temps, devenant
variable explicite), \ill{} :

\begin{itemize}
  \item[g$'$)] \struck{Idem} \uncertain{Analogue} à g), mais avec un « champ de
    vecteurs variable avec le temps » $\xi(t)$ \quad
    ($\underline{X_1 \times X \to \mathbf{T}(X)}$ au-dessus de $X$ via
    $\mathrm{pr}_2$, \struck{formelle \ill{}}), on résolvant l'équation
    différentielle formelle
    \[
      \frac{d}{dt}\, x(t) = \xi(t, x)
    \]
    La solution $x = \varphi(t, x_0)$ correspond\supplied{ant} : un
    $\varphi : X_1 \times X \to X$ qui est un $A$-morphisme.
\end{itemize}

\page{5}

\textbf{Conséquence.} Soit $k_0 = A_0$, c'est un sous-corps de $k$. Soit
$\mathcal{L}(k_0)$ la catégorie des alg.\ de Lie de dim.\ finie sur $k_0$. Par
Hausdorff, on a un foncteur
\[
  \mathcal{L}(k_0) \xrightarrow{\text{Hausdorff}} \mathrm{G\text{-}f}(k_0)
  \xrightarrow{\text{ext.\ de la base}} \mathrm{G\text{-}f}(k)
\]
où « $\mathrm{G\text{-}f}$ » désigne la \textbf{catégorie} des groupes formels.
Ceci dit, ce foncteur se factorise par les $A$-groupes (formels) sur $k$. De
façon plus précise, soit $\overline{W}(\mathfrak{g}_0)$ le complété à l'origine
(de $\mathfrak{g}_0 \otimes_{k_0} k$) \struck{\ill{}} du $k_0$-espace vectoriel
\struck{$\mathfrak{g}_0$}. C'est un $A$-schéma formel.

Les formules de Hausdorff, où $\mathfrak{g}_0$ est une alg.\ de Lie,
définissent une loi de groupe
\[
  \overline{W}(\mathfrak{g}_0) \times \overline{W}(\mathfrak{g}_0) \to
  \overline{W}(\mathfrak{g}_0)
\]
Cette dernière est une $A$-loi ; \struck{\ill{}, il \uncertain{suffit} de \ill{}
\uncertain{des} \ill{} du système de Pfaff} \struck{1°) \ill{} $A$-\ill{}
\uncertain{associé} est donné par une} $A$-loi et si
$\mathfrak{g}_0 \to \mathfrak{h}_0$ est un hom.\ d'algèbres de Lie sur $k_0$,
l'hom.\ de groupes formels
$\overline{W}(\mathfrak{g}_0) \to \overline{W}(\mathfrak{h}_0)$ est un $A$-hom.
Utilisant le th.\ d'Ado, on est ramené à prouver que si
$\mathfrak{g}_0 \subset \mathfrak{gl}(n, k_0)$, alors le s/groupe formel de
$\mathrm{GL}(n, k_0)$ qu'il définit est un sous-$A$-schéma formel….

\page{6}

\textbf{Cor} $A_1 \subset S_1(k)$ stable par primitivation canonique (i.e. celle
qui s'annule à l'origine), contient
$\log(1 + t_1)$, $\exp t_1$, $\cos t_1$, $\sin t_1$, $\mathrm{tg}\, t_1$ ….

Soit $X$ schéma formel lisse de dim.\ $n = p + q$. On sait ce que c'est qu'un
\textbf{feuilletage} (de dim.\ $p$) de $X$ : i.e. variété formelle quotient $B$
de $X$, avec $X \to B$ submersif, $\dim B = q = n - p$. À un tel feuilletage
correspond un sous-module
\[
  \mathfrak{P} \subset \Theta(X)
\]
[champs de vecteurs le long de la fibre] qui est un « système de Pfaff » de dim.\
$p$ complètement intégrable, et inversement. Si $X$ a une $A$-structure, on dit
qu'un feuilletage est un $A$-feuilletage si $\exists$ \uncertain{(loc.)} $B$ une
$A$-structure (d'ailleurs unique, grâce à f)) telle que $X \to B$ soit un
$A$-morphisme. Cela \textbf{implique} que le sous-module
$\mathfrak{P} \subset \Theta(X)$ (ou sous-fibré de $\mathbf{T}(X)$) soit un
$A$-sous-fibré — i.e.\ provient d'un $A$-syst.\ de Pfaff compl.\ intégrable. On
souhaite qu'ait lieu une réciproque, savoir

\marginal{\textbf{NB} h) implique g), qui équivaut à h) pour $p = 1$}

\begin{itemize}
  \item[h)] Si $\mathfrak{P}_A$ est un $A$-syst.\ de Pfaff compl.\ intégrable sur
    le $A$-schéma formel $X$, alors il donne lieu à un $A$-feuilletage.
    \struck{\ill{} réalisé \ill{}}. [En d'autres termes, si $i : Y \to X$ est un
    $A$-morphisme donné tel que le composé
    \[
      \mathbf{T}(Y)_0 \xrightarrow{\mathbf{T}(i)_0} \mathbf{T}(X)_0 \to
      \mathbf{T}(X)_0 / \mathfrak{P}(0)
    \]
    soit un isom., alors l'unique morphisme
    \[
      \pi : X \to Y
    \]
    tel que $\pi \circ i = \mathrm{id}_Y$ et que \uncertain{$\varphi$} soit
    compatible avec $\mathfrak{P}$, est un $A$-morphisme….
\end{itemize}

[$\Leftarrow$ h$'$) la fibre de $0$ dans un tel feuilletage est une
sous-$A$-variété formelle….]

\page{7}

\note{toute la page est barrée de longs traits obliques ; c'est une reprise de la
fin de la page 6, transcrite ici parce qu'elle se lit}

\textbf{Cor} $A_1$ est stable par primitivation canonique (i.e.\ prise
d'augmentation nulle) contient $\log(1 + t_1)$ (primitive de
$\frac{1}{1 + t_1}$) donc $\exp t_1$ (série inverse), les fonctions
$\cos t_1$, $\sin t_1$, $\mathrm{tg}\, t_1$ sont $\in A_1$ ….

\struck{\textbf{Lemme}} \struck{1) Soit} Soit $X$ un $A$-schéma formel (lisse)
\add{de dim.\ $n = p + q$}, $\mathfrak{P}$ un « système de Pfaff de degré $p$
complètement intégrable sur $X$ », \struck{\ill{} dans un $A$-\ill{}}
\struck{Soit $Y$, \ill{} schémas formels (de dim.\ $p$ resp.\ $q$), des
$A$-morphismes, $p : X \to Y$ et $i : Z \to X$ (tels que $\mathbf{T}(p)(0)$ soit
un isom.\ $\mathfrak{P}(0) \to \mathbf{T}(Y)_0$ et que le composé
$\mathbf{T}(Z)_0 \xrightarrow{\mathbf{T}(i)_0} \mathbf{T}(X)_0 \to \mathbf{T}(X)_0 / \mathfrak{P}(0)$
soit un isom. On sait qu'il existe alors un isomorphisme formel (pas néc.
$A$-isom.) unique $X \xrightarrow[\simeq]{\varphi} Y \times Z$
tel que $\mathrm{pr}_1 \circ \varphi = p$, $\varphi \circ i = (z \mapsto (0, z))$)}
\note{la parenthèse qui commence à « Soit $Y$ » est encadrée d'un trait et biffée}

Soit $Z$ \supplied{un $A$-}schéma formel \uncertain{(lisse)} de dim.\ $q$,
$i : Z \to X$ un $A$-morphisme tel que le composé
\[
  \mathbf{T}(Z)_0 \xrightarrow{\mathbf{T}(i)_0} \mathbf{T}(X)_0 \to
  \mathbf{T}(X)_0 / \mathfrak{P}(0)
\]
soit un isom. On sait qu'il existe alors un unique morphisme formel
\[
  X \xrightarrow{\pi} Z
\]
tel que $\pi \circ i = \mathrm{id}$, et que « $\pi$ soit \struck{constant sur les}
compatible avec le feuilletage défini par $\mathfrak{P}$ ». Conditions
équivalentes

\begin{itemize}
  \item[(i)] $\pi$ est un $A$-morphisme
  \item[(ii)] $\exists$ $A$-isom.\ $X \simeq B \times Y$, tel que le feuilletage
    formel de $X$ défini par $\mathfrak{P}$ soit le feuilletage défini par la
    projection $\mathrm{pr}_1 : X \to B$. On dit
\end{itemize}


\section{Reprise : $k$ anneau commutatif, schémas $A_*$-formels}

\note{ce titre est de l'édition}

\page{9}

\note{nouvelle rédaction, reprenant la page 2 avec $k$ un anneau commutatif
quelconque}

$k$ anneau commutatif

$S_n = k[[t_1, \ldots, t_n]]$ \quad ($n \geq 0$)

$X_n = \mathrm{Spf}(S_n)$

$X_* = (X_n)_{n \geq 0}$ \quad schéma formel simplicial \textbf{augmenté}

$S_* = (S_n)_{n \geq 0}$ \quad anneau top.\ cosimplicial coaugmenté

On suppose donné, pour tout $n \geq 0$, un sous-ensemble
\[
  A_n \subset S_n
\]
On pose $A_* = (A_n)_{n \geq 0}$. Un morphisme $X_m \xrightarrow{p} X_n$,
associé à un hom.\ de $k$-algèbres topologiques $S_n \xrightarrow{p^*} S_m$, est
dit un $A_*$-\textbf{morphisme} si
\[
  p^*(t_i) = f_i \in S_m \text{ est } \in A_m \text{ pour tout } 1 \leq i \leq n .
\]
Donc cela signifie que les $\mathrm{pr}_i \circ p$ ($1 \leq i \leq n$) sont des
$A_*$-morph. On suppose :

\marginal{\textbf{NB} les $f_i$ \ill{} les $\varepsilon(f_i)$
\uncertain{arbitraires}}

\begin{itemize}
  \item[(a)] \marginal{Identités} Les opérations simpliciales entre les $X_n$
    sont des $A_*$-morphismes ($\Leftrightarrow$ $\forall n \geq 0$ et
    $1 \leq i \leq n$, $t_i \in A_n$ ; $\Leftrightarrow$ les morphismes
    identiques des $X_n$ sont $\in A_n$)
  \item[(b)] \marginal{Composition} Le composé de deux $A_*$-morphismes est un
    $A_*$-morphisme [$\Leftrightarrow$ si $p : X_m \to X_n$ est un
    $A_*$-morphisme, alors $p^*(A_n) \subset A_m$ $\Leftrightarrow$ si
    $f_1, \ldots, f_n \in A_m$, et $\varepsilon(f_1), \ldots, \varepsilon(f_n)$
    nilpotents \add{avec $\varepsilon$ l'augmentation}, alors pour tout
    $\varphi \in A_n$, on a $\varphi(f_1, \ldots, f_n) \in A_m$]
    \marginal{\textbf{NB} \struck{\ill{}} Ce serait vrai (sans \uncertain{supposer}
    $\varepsilon(\varphi)$ nilpotent) \uncertain{dans} \ill{}
    \struck{si $a \in A_0$}, sachant que
    $\mathrm{Im}(A_0 \to S_n) \subset A_n$}
  \item[(c)] \struck{$1 \in A_0$, et} Le morphisme « différence »
    $X_2 = X_1 \times_k X_1 \to X_1$ ($(x, y) \mapsto x - y$) et « produit »
    [$(x, y) \mapsto xy$] sont des $A_*$-morphismes
    [$\Leftrightarrow$ $t_1 - t_2$ et $t_1 t_2$ sont $\in A_2$]
    \add{et $1 \in A_0$} $\Leftrightarrow$ les $A_n$ sont \struck{$A$-\ill{}}
    des sous-anneaux unitaires des $S_n$
    \struck{\ill{} stables par produits, et par différences, \ill{} \ill{}
    $1 \in S_0$}
    \marginal{condition \uncertain{anodine}}
\end{itemize}

\note{l'énoncé (c) est encadré et traversé de traits obliques, sans que rien le
remplace ; la page 10 le suppose acquis}

\page{10}

\textbf{NB} \struck{Le plus petit $A_*$} a) b) c) impliquent que $A_*$ est un
sous-anneau cosimplicial de $S_*$. Le plus petit $A_*$ satisfaisant a) b) c) est
formé des
\[
  A_n = \mathrm{Im}\bigl(\mathbf{Z}[t_1, \ldots, t_n] \to k[[t_1, \ldots, t_n]]\bigr)
  \simeq \Lambda[t_1, \ldots, t_n]
\]
où
\[
  \Lambda = \mathrm{Im}(\mathbf{Z} \to k)
  \begin{cases}
    \simeq \mathbf{Z} & \text{ou} \\
    \simeq \mathbf{Z}/n\mathbf{Z} & \text{pour } n \in \mathbf{N}^* \text{ convenable.}
  \end{cases}
\]

Avec les conditions a) b) c) on construit la catégorie
$\mathrm{Sfl}_{A_*}$ des $A_*$-schémas formels \uncertain{lisses}. Dans cette
catégorie, il y a des produits finis ($X_n \simeq (X_1)^n$), et des produits
fibrés, dans le cas « transversal ».

\marginal{Le foncteur « anneau des $A_*$ » commute aux limites $\leftarrow$}

On obtient un foncteur $\mathbb{V}_{A_*}$ de la cat.\ des $k_0$-modules libres de
t.f.\ \struck{est} (où $k_0 = A_0 \subset S_0 = k$) \uncertain{(\ill{} ?)} vers
les $A_*$-schémas formels, par $M_0 \mapsto \mathrm{Spf}\,
\widehat{\mathrm{Sym}}{}^*_k(M_0 \otimes k)$ (i.e.\
$M_k = M_0 \otimes_{k_0} k$). C'est un foncteur \uncertain{pleinement fidèle}
\ill{} qui transforme les objets groupes de $\mathrm{Lib}(k_0)$ en les groupes
$A_*$-formels.

\marginal{\struck{\ill{} \ill{} \ill{} schémas $A_*$ \ill{} \ill{} plus gén.\ \ill{}}}
\note{la note encadrée en biais au milieu de la marge gauche est rayée de
traits verticaux ; seuls quelques mots se lisent}

Plus gén., si $\mathrm{Lin}(k_0)$ est la catégorie des $k_0$-schémas
\textbf{augmentés} (\uncertain{\ill{}} des $\mathrm{Spf}\,
\widehat{\mathrm{Sym}}{}^*_k(M_0)$, \ill{}), il y a un foncteur
$X \to \widehat{X}^{(A_*)}$ de $\mathrm{Lin}(k_0)$ vers les $A_*$-schémas
formels, commutant aux produits finis, donc transformant objets groupes en
objets groupes \ldots

\marginal{$\mathcal{O}_{A_*}(X) \subset \mathcal{O}(X)$, où
$\mathcal{O}_{A_*}(X) = \mathrm{Hom}_{A_*}(X, X_1)$ et
$\mathcal{O}(X) = \mathrm{Hom}_{k\text{-formel}}(X, X_1)$\note{les deux égalités
sont écrites par des « $=$ » verticaux sous chaque membre} ;
\struck{\ill{}} $A_*$-formel \uncertain{lisse} \ill{} ; Modules $A_*$-formels
sur les schémas $A_*$-formels ; \struck{Le \ill{}} fibré vectoriel $A_*$-formel
$\mathbb{V}_{A_*}(M) := W_{A_*}(M)$}
\note{cette dernière note court verticalement le long de la marge gauche, dans
le bas de la page}

\page{11}

Les conditions a) b) c) sont les \textbf{conditions} \struck{algébriques}
élémentaires. On veut de plus

\marginal{\uncertain{avec} d) \uncertain{condition de coefficientabilité}}

\begin{itemize}
  \item[d)] (Condition hensélienne) Si $u : X_n \to X_n$ est un $A_*$-morphisme
    qui est un isomorphisme formel (i.e.\ si $f_1, \ldots, f_n \in A_n$ sont tels
    que $\varepsilon(f_1), \ldots, \varepsilon(f_n)$ soient nilpotents dans $k$ et
    \[
      \varepsilon\Bigl(\det\Bigl(\frac{\partial f_i}{\partial t_j}\Bigr)\Bigr)
      = \det\Bigl(\varepsilon\Bigl(\frac{\partial f_i}{\partial t_j}\Bigr)\Bigr)
      = \det\Bigl(\frac{\partial f_i}{\partial t_j}(0)\Bigr)
    \]
    est un élément inv.\ de $k$) alors $u^{-1}$ est aussi un $A_*$-morphisme
    ($\Leftrightarrow$ le foncteur des $A_*$-schémas $A_*$-formels vers les
    schémas formels sur $k$ est conservatif).
\end{itemize}

\textbf{Cor.\ 1} Si $f \in A_n$ et $f$ inversible dans $S_n$, i.e.\
$\varepsilon(f) = f(0)$ inversible dans $k$, alors $f^{-1} \in A_n$ ;
$\Leftrightarrow$ \struck{\ill{} \ill{} \ill{}} si $f \in A_n$,
$\varepsilon(f) = 0$, alors $(1 + f)^{-1} \in A_n$ [$\Leftrightarrow$ si
$f \in A_n$ et $\varepsilon(f) = 1$, alors $f^{-1} \in A_n$]
$\Leftrightarrow$ $A_n \cap \mathfrak{m}_n \subset \mathrm{rad}\, A_n$, où
\struck{\ill{}} $\mathfrak{m}_n = (t_1, \ldots, t_n) S_n$ \struck{\ill{}}]

\marginal{moyennant condition d) (coefficientabilité)}

En effet, considérons la \struck{multi\ill{}} \supplied{l'}automorphisme
\[
  u : (x_1, \ldots, x_n, x_{n+1}) \mapsto
  (x_1, \ldots, x_n, f(x_1, \ldots, x_n) . x_{n+1})
\]
de $X_{n+1}$, il est un isom.\ formel d'inverse
\[
  (x_1, \ldots, x_n, x_{n+1}) \mapsto
  (x_1, \ldots, x_n, f(x_1, \ldots, x_n)^{-1} x_{n+1}),
\]
donc comme \struck{\ill{}} $u$ est un $A$-morphisme, par d) il en est de même
de $u^{-1}$, donc
$g(t_1, \ldots, t_{n+1}) = f(t_1, \ldots, t_n)^{-1} t_{n+1} \in A_{n+1}$,
\struck{i.e.\ $X_{n+1} \to X_1$ \ill{}} donc $g(t_1, \ldots, t_n, 1) \in A_n$
i.e.\ $f^{-1} \in A_n$, cqfd.

\marginal{\uncertain{\ill{}} de la condition de coefficientabilité $\uparrow$ :
inconvénient, $1$ n'est pas nilpotent ! Il faut exiger une propriété de
divisibilité : $t_{n+1} f \in A_{n+1} \Rightarrow f \in A_n$}
\note{la note encadrée en bas de la marge gauche, reliée par une flèche au
passage « donc $g(t_1, \ldots, t_n, 1) \in A_n$ », relève que $1$ n'est pas
nilpotent, de sorte que la substitution $t_{n+1} = 1$ n'est pas permise par b)}

\textbf{Cor.\ 2} Soit $P(Z) = Z^N + f_1 Z^{N-1} + \cdots + f_N$,
$\in A_n[Z]$ polynôme unitaire, soit $\zeta \in k_0$ tel que
\[
  P^\varepsilon(\zeta) = \zeta^N + a_1 \zeta^{N-1} + \cdots + a_N = 0 ,
  \quad \text{où } a_i = \varepsilon(f_i) = f_i(0) \in k_0 ,
\]

\page{12}

et $(P^\varepsilon)'(\zeta) = N \zeta^{N-1} + (N-1) a_1 \zeta^{N-2} + \cdots
+ a_{N-1}$, soit $\in k^*$ (donc en fait $\in k_0^*$), de sorte qu'il existe un
unique $F \in S_n$ telle que $P(F) = 0$ et $\varepsilon(F) = \zeta$. Sous ces
conditions, on a $F \in A_n$. [En d'autres termes,
$(A_n, \mathfrak{p}_n = \mathfrak{m}_n \cap A_n)$ est (non seulement
zariskien, selon le cor.\ 1, mais même) hensélien]

\struck{\ill{} \ill{} $Z$ par $Z - \zeta_0$, sps $\zeta_0 = 0$ donc
$f_N = \varepsilon(f_N) = 0$. En effet, on regarde \ill{}
$X_{n+1} \to X_{n+1}$ donné par
$(x_1, \ldots, x_n, x_{n+1}) \mapsto (x_1, \ldots, x_n, x_{n+1}^{N+1} + f_1(x_1, \ldots, x_n) x_{n+1} \ldots)$.
L'hyp.\ sur}
\note{passage encadré d'un trait en escalier et biffé}

\textbf{Plus généralement}

\textbf{Cor.\ 3} Soit $P \in A_{n+1}[t_1, \ldots, t_n, Z]$,
$P = f_0 + f_1 Z + f_2 Z^2 + \cdots$ [les $f_i \in S_n(t_1, \ldots, t_n)$]
supposons $\varepsilon(f_0) = 0$ i.e.\ $\varepsilon(P) = 0$, et
$\varepsilon(f_1) \in k^*$ i.e.\
$\varepsilon\bigl(\frac{\partial P}{\partial Z}\bigr) \in k^*$, de sorte qu'il existe
une unique $F \in S_n$ telle que $P(t_1, \ldots, t_n, F(t_1, \ldots, t_n)) = 0$
et $\varepsilon(F) = 0$. Alors $F \in A_n$.

\marginal{a) b) c) d) impliquent-ils que $S_n$ est fid.\ plat \struck{\ill{}} sur
$A_n$ (si $k$ \uncertain{\ill{}} \uncertain{fini} sur $k_0$) ?}

Soit en effet $u : X_{n+1} \to X_{n+1}$, donné par
$(x_1, \ldots, x_n, \zeta) \mapsto (x_1, \ldots, x_n, P(x_1, \ldots, x_n, \zeta))$,
alors l'hyp.\ dit que $u$ est compatible avec les projections, et que $u$ est un
isomorphisme formel. Comme l'hyp.\ $P \in A_{n+1}$ implique $u$ un
$A_*$-morphisme, $u^{-1}$ est un $A_*$-morphisme, et par suite $u^{-1}$ est un
$A_*$-morphisme. Donc
$u^{-1}(x_1, \ldots, x_n, \zeta) = (x_1, \ldots, x_n, \varphi(x_1, \ldots, x_n, \zeta))$
avec $\varphi \in A_{n+1}$, $(t_1, \ldots, t_n, t_{n+1})$, donc
$F = \varphi(t_1, \ldots, t_n, 0) \in A_n(t_1, \ldots, t_n)$.

\struck{d) (Coefficients différentiables) les $A_n$ sont stables par les
dérivations différentielles standard
$D_{p_j} = D(p_1, \ldots, p_n)$
$\bigl(= \frac{1}{p_1! \cdots p_n!} \frac{\partial^{|p|}}{\partial x_1^{p_1} \cdots \partial x_n^{p_n}}\bigr)$}
\note{cet énoncé, en bas de page, est barré de traits obliques}

\page{13}

Soient $X$ un schéma $A_*$-formel et $Y$ un sous-schéma formel (lisse) de $X$.
\struck{On dit} On dit que $Y$ est un \struck{sous} $A_*$-schéma $A_*$-formel s'il
existe sur $Y$ une $A_*$-structure telle que l'inclusion canonique soit un
$A_*$-morphisme. Cette $A_*$-structure est alors unique. En effet, il est
immédiat qu'on peut (sans hyp.\ sur $Y$) trouver un $A_*$-morphisme
$X \to Y'$ ($Y'$ un schéma $A_*$-formel) tel que le morphisme induit
$Y \to Y'$ soit un isomorphisme formel. \struck{Considérons} Considérons
l'isomorphisme inverse, et le composé $Y' \to Y \to X$. Utilisant
\uncertain{(d)}, on voit que \struck{\ill{}} $Y$ est un sous-schéma $A_*$-formel
ssi $Y' \to X$ est un $A_*$-morphisme, et alors la $A_*$-structure de $Y$ est
celle déduite.

\textbf{Exemples} La diagonale de $X \times X$ est un sous-$A_*$-schéma formel.
Plus gén.\ celle de $X \times_Y X$ si $X \to Y$ est un $A_*$-morphisme submersif
\struck{(\ill{} formel \ill{})}.

\marginal{Un morphisme formel $X \to Y$ entre schémas $A_*$-formels est un
$A_*$-morphisme ssi son graphe est un sous-schéma $A_*$-formel de $X \times Y$}

Soit $\mathcal{X}$ un schéma sur $\mathcal{O}_{A_*}(X)$ ($X$ un $A_*$-schéma
formel) muni d'une section $s_0$ sur $\mathcal{O}_{A_*}(X)$, supposons que
$\mathcal{X}$ soit lisse sur $\mathrm{Spec}\, \mathcal{O}_{A_*}(X)$ aux points de
cette section. Considérons le complété formel $\widehat{\mathcal{X}}$ de
$\mathcal{X} \otimes_{\mathcal{O}_{A_*}(X)} \mathcal{O}(X)$ le long de la section
$s$ déduite de $s_0$. Alors une construction utilisant \uncertain{e)} permet
d'y mettre une $A_*$-structure, telle que
$X \to \widehat{\mathcal{X}} \to X$ soient des $A_*$-morphismes.
\struck{\ill{}} $\widehat{\mathcal{X}}$ dépend $A_*$-fonctoriellement des triples
$(X, \mathcal{X}, s_0)$, sa formation est compatible avec les images inverses
par $X' \to X$ \ldots)
\note{la phrase « telle que $X \to \widehat{\mathcal{X}} \to X$ soient des
$A_*$-morphismes » est ajoutée dans l'interligne}

\page{14}

\begin{itemize}
  \item[\uncertain{d)}] (Conditions différentielles) Tout $A_n$ est stable par
    les op.\ différentiels élémentaires $D_p$
    [$\Leftrightarrow$ si $f(x_1, \ldots, x_n, y_1, \ldots, y_m) \in A_{n+m}$ est
    mis sous la forme $\sum_{p \in \mathbf{N}^m} f_p(x)\, y^p$, alors les $f_p$
    sont $\in A_n$]
\end{itemize}
\note{la lettre de cet énoncé se lit « d) », qui est déjà celle de la
condition hensélienne (page 11) ; la page 13 invoque « (d) » puis « e) », et la
page 15 enchaîne sur « f) (Condition intégrale) » : c'est vraisemblablement
l'énoncé e) de la série}

\marginal{\uncertain{les conjugués} \ill{} \uncertain{les coefficients}
\ill{} ?}

Ceci implique que les algèbres $\mathcal{P}^n_{X/k}$ sur $\mathcal{O}(X)$ ($X$ un
schéma $A_*$-formel) ont \struck{une} des $A_*$-structures canoniques et
fonctorielles, de même $\mathcal{P}^n_{X/Y}$ si $X \to Y$ est un
$A_*$-morphisme submersif, de même \struck{\ill{}} si $X \supset Y$ est un
couple de schémas $A_*$-formels, $\mathcal{O}(Y) = \mathcal{O}(X)/\mathcal{J}$,
les invariants normaux $\mathcal{O}_X/\mathcal{J}^n$ ont des $A_*$-structures
canoniques, fonctorielles, compatibles avec images inverses transversales. Les
op.\ diff.\ universel et relatif
$\mathcal{O}(X) \to \mathcal{P}^n_{X/k}$ etc.\ \struck{\ill{} $A_*$-morphismes},
appliquent $\mathcal{O}_{A_*}(X)$ dans $(\mathcal{P}^n_{X/k})_{A_*}$ \ldots

\uncertain{Si} les $\mathrm{Diff}^n(X/k)$ (duaux des $\mathcal{P}^n_{X/k}$) ont
\struck{\ill{}} des $A_*$-structures canoniques, leur composition
\struck{est} laisse stables les parties $\mathrm{Diff}^n(X_{A_*}/k)$ des
$A_*$-op.\ diff., et ces derniers stabilisent \struck{les $A_*$ \ill{}}
$\mathcal{O}_{A_*}(X)$ \ldots\ Item pour les op.\ diff.\ relatifs, pour
$X \to Y$ $A_*$-morphisme \struck{\ill{}} submersif.

\textbf{Cor} $A_n \subset k_0[[t_1, \ldots, t_n]]$. \struck{\ill{}} D'où un
foncteur allant \struck{\ill{}} des schémas $A_*$-formels vers les schémas
formels sur $k_0$ \ldots
\struck{Soit $X \to Y$ un morphisme submersif de schémas formels, \ill{} $X$}

\page{15}

Soit $p : X \to Y$ un morphisme submersif de schémas formels (lisses),
supposons $X$ muni d'une $A_*$-structure. On dit que $Y$ admet une
$A_*$-structure quotient de celle de $X$ (moyennant $p$) s'il admet une
$A_*$-structure telle que $p$ devienne un $A_*$-morphisme. Celle-ci est alors
unique : en effet, il existe (sans condition sur $Y$) un \struck{\ill{}}
sous-schéma $A_*$-formel $Y'$ de $X$, tel que le morphisme formel induit
$Y' \to Y$ soit un isomorphisme. \uncertain{Considérons} le composé
$X \xrightarrow{p} Y \xrightarrow{\sim} Y'$, on voit alors grâce à
\uncertain{e)} que la $A_*$-structure quotient existe ssi $X \to Y'$ est un
$A_*$-morphisme, et la $A_*$-structure sur $Y$ est celle déduite de celle de
$Y'$.
\note{« Considérons » est ajouté dans l'interligne, et le composé est écrit
$X \to Y'$ avec $p$, $Y$ et une flèche portés au-dessus ; un mot après $Y'$ est
biffé}

\struck{\ill{}} Soit $X \xrightarrow{p} Y$ un morphisme formel submersif. Alors
$R = X \mathbin{\widehat{\times}_Y} X \subset X \mathbin{\widehat{\times}} X$
est un sous-schéma formel lisse, et c'est une relation d'équivalence [stable
par symétrie, contenant la diagonale, transitif\ldots]. \struck{Inversement}
Plus précisément, $X$ étant donné, il y a corr.\ biunivoque entre ces quotients
formels $Y$, et les $R$ rel.\ d'équiv.\ formelles. Si $X$ est muni d'une
$A_*$-structure, alors si $Y$ admet une $A_*$-structure quotient, on voit que
$R$ est un sous-schéma $A_*$-formel de $X \mathbin{\widehat{\times}} X$ i.e.\ il
a une $A_*$-structure induite.

\begin{itemize}
  \item[f)] (Condition \textbf{intégrale}) La réciproque est vraie, i.e.\ si
    $R$ est une relation d'équival.\ formelle dans $X$, qui soit une
    $A_*$-sous-structure de $X \mathbin{\widehat{\times}} X$, alors
\end{itemize}

\page{16}

le quotient formel $Y = X/R$ admet une $A_*$-structure quotient.

[Heuristiquement, étant donné un $A_*$-feuilletage de $X$, la variété formelle
des feuilles admet une $A_*$-structure quotient. En car.\ $0$, il est « bien
connu » que les $A_*$-feuilletages correspondent aux « systèmes de Pfaff
complètement \textbf{intégrables} » sur $X$, et la « condition intégrale »
\textbf{signifie} que l'\textbf{intégration} de ce système (effectuée, disons,
par le morphisme formel plus haut $X \to Y'$) se fait par un
$A_*$-morphisme\ldots]
\note{la page s'arrête ici ; la moitié inférieure du feuillet est blanche}


\section{Théorème de préparation « formel »}

\note{ce titre est de lui : « Th. de préparation "formel" », souligné, ouvre la
page 18}

\page{18}

\textbf{Th.\ de préparation « formel »} Soit $A$ un anneau linéairement
topologique, séparé et complet, $Z$ une indéterminée,
$F(Z) = \sum_{i \geq 0} a_i Z^i \in A[[Z]]$, $N \in \mathbf{N}$ tel que

\begin{itemize}
  \item[(i)] $a_i$ est top.\ nilpotent pour $0 \leq i \leq N-1$
  \item[(ii)] $a_N$ est inversible
\end{itemize}

Alors $F$ est un élément \textbf{régulier} de $A[[Z]]$ et $A[[Z]]/F$ a comme base
$(Z^i)_{0 \leq i \leq N-1}$, en d'autres termes tout élément $G$ de $A[[Z]]$
s'écrit de façon unique sous la forme
\[
  G = QF + \sum_{0 \leq i \leq N-1} b_i Z^i \qquad
  Q \in A[[Z]],\ b_i \in A
\]
\note{l'énoncé et sa conclusion sont marqués chacun d'un trait vertical dans la
marge gauche}

\textbf{Démonstration} Comme $A$ est séparé et complet pour sa top.\ linéaire,
\uncertain{on est ramené au} cas où $A$ est discret, alors (i)
\textbf{signifie} que les $a_i$ ($0 \leq i \leq N-1$) sont nilpotents. Soit
$\mathfrak{m}$ l'idéal qu'ils engendrent, $A_0 = A/\mathfrak{m}$, et désignons
par $\dot{a}_i$, $\dot{F}$ etc.\ les réductions mod $\mathfrak{m}$.
\struck{\ill{}} Donc
\[
  \dot{a}_i = 0 \quad \text{pour } 0 \leq i \leq N-1
\]
i.e.
\[
  \dot{F} = \dot{a}_N Z^N + \cdots = \dot{a}_N Z^N (1 + \cdots)
\]
donc, comme $\dot{a}_N$ est inversible par (ii), on voit que $\dot{F}$ est de la
forme $\dot{F} = \dot{u} Z^N$.
\note{les indices de $\dot{a}_N$ et les exposants de $Z^N$ sont surchargés : un
« $N+1$ » semble corrigé en « $N$ »}

\struck{\ill{}} Donc
\[
  A_0[[Z]]/\dot{F} \simeq A_0[[Z]]/Z^N
\]
\struck{$\simeq A_0[Z]/Z^N$}
a comme base les $Z^i$ ($0 \leq i \leq N-1$).

\struck{i.e. (i) les $Z^i$ ($0 \leq i \leq N-1$) engendrent}
\struck{$A[[Z]]/F A[[Z]]$.}
\struck{En effet, ils engendrent mod $\mathfrak{m}$,}
\struck{on gagne donc par Nakayama.} \struck{\ill{}}
\struck{(ii) $F$ est régulier. En effet,} \struck{\ill{} associé \ill{}}
\struck{$\mathrm{Gr}^*(A)^N$} \struck{\ill{}}
\struck{Soit $Q \in A[[Z]]/\mathfrak{m}$} \struck{\ill{}}
\struck{$F$ divise} \struck{\ill{}}
\struck{quand, correspond : l'op.\ de} \struck{\ill{}}
\struck{$F$, qui est injective OK.}
\note{tout le bas de la page, à partir de « (i) les $Z^i$ », est enfermé dans
un cadre, traversé de lignes horizontales et barré de grandes diagonales
croisées ; seuls des fragments se lisent}

\page{19}

\struck{(iii) Étude du cas où $a_i = 0$ pour $i \geq N+1$, $a_N = 1$,}
\struck{i.e.\ dans ce cas} \struck{$F = Z^N + a_{N-1} Z^{N-1}$} \struck{$+ \cdots + a_0$,}
\struck{les $a_i$ ($0 \leq i \leq N-1$) nilpotents.}
\struck{$A[Z]/F A[Z] \to A[[Z]]/F A[[Z]]$,}
\struck{base $(Z^i)_{0 \leq i \leq N-1}$ sur $A$,}
\struck{est donc surjectif}
\struck{(car les $Z^i$ ($0 \leq i \leq N-1$) engendrent}
\struck{$A[[Z]]/F A[[Z]]$ comme $A$-module) par (i).}
\struck{Pour qu'il soit bijectif il} \uncertain{faut et il suffit}
\struck{que $(F A[[Z]]) \cap A[Z] = F A[Z]$}
\struck{(il revient au même de dire que les} \struck{\ill{}}
\struck{$Z^i$ ($0 \leq i \leq N-1$) forment une base}
\struck{de $A[[Z]]/F A[[Z]]$).}
\struck{Soit donc $a_0 + a_1 Z + \cdots + a_{N-1} Z^{N-1}$}
\struck{qui soit multiple de $F$ dans $A[[Z]]$ i.e.}
\note{le haut de la page, jusqu'à « multiple de $F$ dans $A[[Z]]$ », est
barré de deux longues obliques}

Considérons l'application
\[
  \begin{array}{rccc}
    \varphi : & A^N + A[[Z]] & \longrightarrow & A[[Z]] \\
    & \bigl((b_i)_{0 \leq i \leq N-1}, Q\bigr) & \longmapsto &
      \sum_{0 \leq i \leq N-1} b_i Z^i + QF
  \end{array}
\]
l'application déduite par passage aux gradués associés pour la filtration
$\mathfrak{m}$-adique (finie) est
\[
  (\mathrm{Gr}\, A)^N \times (\mathrm{Gr}\, A)[[Z]] \to (\mathrm{Gr}\, A)[[Z]]
\]
donnée par un degré $n$ par
\[
  \bigl((b_i), Q\bigr) \longmapsto \sum_{0 \leq i \leq N-1} b_i Z^i + Q \dot{F}
\]
et comme $\dot{F} = \dot{a}_N \alpha Z^N$ ($\alpha$ \ill{} $\in A_0^*$)\note{la
parenthèse, surchargée, porte après $\alpha$ un signe et un « $a_{N+1}$ » ou
« $a_N$ » corrigé ; ce que vaut $\alpha$ ne se lit pas}, ceci est manifestement bijectif,
donc $\varphi$ aussi, cqfd.

\end{document}
